A complete \texttt{bisect} redistricting run---census data fetch, graph construction, recursive bisection, analysis, and report generation---costs under \$50,000 total for a typical state, including staff time, compute, and independent audit. Census tract boundary data (40GB+ download) is available free from the Census Bureau. A 50-state congressional redistricting run completes in approximately 18 minutes of wall-clock compute time on a modern 8-core machine, at a cloud cost of under \$200. The dominant cost is staff time for quality assurance, public engagement, and legal review---not the algorithm itself. We compare these costs against two benchmarks: Iowa's Legislative Services Agency model (\$3 million per cycle), which is the current gold standard for low-cost non-partisan redistricting, and California's Citizens Redistricting Commission (\$35 million+ per cycle), which is the highest-cost commission model. Algorithmic redistricting is 60--700 times cheaper than existing models. We analyze how the \texttt{bisect} manifest format satisfies state and federal administrative procedure act (APA) record-keeping requirements, including the factual record standard of \textit{Motor Vehicle Manufacturers Association v. State Farm}, 463 U.S.\ 29 (1983). We provide a model procurement specification that state redistricting offices can use to contract for algorithmic redistricting services or to evaluate vendor proposals.
